\documentclass[conference]{IEEEtran}

\usepackage{amsmath}
\usepackage{amssymb}
\usepackage{booktabs}
\usepackage{multirow}
\usepackage{url}
\usepackage{algorithm}
\usepackage{algorithmic}
\usepackage{graphicx}
\usepackage{microtype}

\begin{document}

\title{Zero-Shot Prompt Injection Detection via N-gram Pooling of Token-Level Hidden States}

\author{
  \IEEEauthorblockN{Lokesh M}
  \IEEEauthorblockA{TODO: Affiliation}
}

\maketitle

\begin{abstract}
Retrieval-augmented generation (RAG) pipelines are vulnerable to prompt injection
attacks, where adversaries embed malicious instructions inside retrieved documents.
Existing embedding-based detectors rely on mean-pooled sentence vectors, which
suppress short injection spans hidden within long benign text --- a phenomenon we
formalize as \emph{semantic dilution}. We extract per-token hidden states from a
frozen 22M-parameter sentence transformer already present in the pipeline and apply
N-gram pooling over sliding windows to localize injection signals. On a controlled
chunk-level benchmark (1,575 chunks, Enron corpus), our method achieves
AUC\,=\,0.901 (95\,\% CI [0.869, 0.930]), with confidence intervals non-overlapping
with the mean-pool baseline (AUC\,=\,0.808). On the public xTRam1 dataset, zero-shot
AUC\,=\,0.870 exceeds the supervised baseline of 0.764. All three scoring variants
execute at identical CPU latency (\(\approx\)44\,ms per 500-character chunk). The
method requires no labeled data, no fine-tuning, and no access to the generative LLM
at inference time.
\end{abstract}

\section{Introduction}
\label{sec:intro}

Retrieval-augmented generation augments language model responses by injecting
retrieved document passages into the context window. This design surface is
exploitable: an adversary who can write to any document in the knowledge base can
embed a prompt injection that hijacks the model's behavior when that document is
retrieved~\cite{yi2025bipia}. Defending against such attacks at retrieval time --- before the
generative LLM is ever invoked --- is therefore a critical pipeline hardening step.

Embedding-based detectors compare a semantic representation of each retrieved chunk
against a small set of known-malicious reference phrases~\cite{ayub2024embedding}. The
appeal is clear: the sentence transformer used for retrieval is already in the
pipeline, so detection adds negligible infrastructure. However, all prior
embedding-based work collapses each chunk to a single mean-pooled
vector~\cite{ayub2024embedding,alamsabi2026context}. When a short injection phrase is
surrounded by a long benign passage, the mean-pooled vector converges toward the
benign centroid and the injection signal is washed out. We call this \emph{semantic
dilution} and formalize it in Section~\ref{sec:method}.

Methods that exploit sub-sentence granularity do exist. InstructDetector~\cite{wen2025instructdetector}
extracts per-token hidden states and gradients from the guarded generative LLM and
achieves near-perfect accuracy. But this design requires a full forward-and-backward
pass through the LLM at inference time --- expensive, and unavailable at ingestion time
before any query is issued.

We propose a middle path. A sentence transformer already produces contextual per-token
hidden states as a byproduct of its forward pass. Rather than collapsing these to a
sentence vector, we apply N-gram pooling: we average each window of \(N\) consecutive
token embeddings and score the chunk against a reference set using the maximum
similarity over all windows. This recovers phrase-level intent that mean-pooling
destroys, while remaining entirely within the lightweight encoder.

Our contributions are:
\begin{enumerate}
  \item A formal characterization of semantic dilution in mean-pooled injection
        detection (Section~\ref{sec:method}).
  \item N-gram pooling over token-level hidden states for localized, zero-shot injection
        scoring, with no training on injection data (Section~\ref{sec:method}).
  \item A controlled chunk-level evaluation with bootstrap confidence intervals on the
        Enron corpus, demonstrating statistically significant improvement over mean-pool
        and per-token baselines (Section~\ref{sec:results}).
  \item A zero-shot comparison against supervised baselines on three public benchmarks,
        with the N-gram scanner exceeding the supervised AUC on xTRam1
        (Section~\ref{sec:results}).
\end{enumerate}

\section{Related Work}
\label{sec:related}

\subsection{Embedding-Based Injection Detection}

Ayub and Majumdar~\cite{ayub2024embedding} apply the same encoder used in this work
(\texttt{all-MiniLM-L6-v2}) to injection detection by computing mean-pooled sentence
embeddings and training a random forest classifier on top. Their best supervised AUC
is 0.764. The key limitation is that mean pooling is applied before classification,
discarding positional structure within the chunk.

Alamsabi et al.~\cite{alamsabi2026context} construct context-pair embeddings from
consecutive paragraph segments and train a classifier on the BIPIA
benchmark~\cite{yi2025bipia}, achieving F1\,=\,0.977. Their method is trained on the
same distribution it is evaluated on; it requires labeled examples of injections per
domain and cannot generalize zero-shot to unseen injection styles.

\subsection{Token-Level Hidden State Methods}

InstructDetector~\cite{wen2025instructdetector} is the closest prior work exploiting
sub-sentence granularity. It extracts per-token hidden states and gradient magnitudes
from Llama-3.1-8B and trains a classifier on top, achieving 99.6\,\% accuracy.
The critical constraint is the requirement for a forward-and-backward pass through the
guarded generative LLM: this is unavailable at document ingestion time, and adds
inference cost proportional to LLM scale. Our approach uses a 22M-parameter encoder
and requires no gradients.

\subsection{RAG Security Benchmarks}

Yi et al.~\cite{yi2025bipia} introduce BIPIA, the standard benchmark for indirect
prompt injection in RAG contexts, which provides 10,000 attack-response pairs across
five tasks. The JailbreakBench framework~\cite{chao2024jailbreakbench} defines 100
harmful behaviors across 10 taxonomic categories; we use it as the source of our
reference set and injection corpus. The sentence-transformers
library~\cite{reimers2019sentence} provides the \texttt{all-MiniLM-L6-v2} encoder
used throughout this work.

\section{Method}
\label{sec:method}

\subsection{Problem Statement}

Let \(D\) be a document retrieved from an external corpus during a RAG pipeline
query. \(D\) is partitioned into \(K\) non-overlapping chunks
\(c_1, \dots, c_K\) of at most \(W\) characters. A chunk \(c_k\) is
\emph{poisoned} if an adversary has embedded an injection phrase within the
surrounding benign text; otherwise it is \emph{clean}. Our goal is to assign a
scalar score \(s(c_k) \in [0, 1]\) such that poisoned chunks score higher than
clean ones --- without any training on labeled injection data and without access to
the generative LLM.

\subsection{Token-Level Hidden States}

We use a lightweight sentence transformer \(f\) (\texttt{all-MiniLM-L6-v2}, 22M
parameters, \(d = 384\)) already present in the RAG pipeline for retrieval. For an
input chunk \(c\) tokenized into \(L\) subword tokens, \(f\) produces contextual
token embeddings:
\[
  H = f(c) \in \mathbb{R}^{L \times d}
\]
where \(h_i \in \mathbb{R}^d\) is the hidden state of the \(i\)-th token.
The model is used in frozen inference mode; no gradient computation is performed.

\subsection{Semantic Dilution}

Standard sentence-level scoring computes a mean-pooled representation:
\[
  \mathbf{g} = \mathrm{normalize}\!\left( \frac{1}{L} \sum_{i=1}^{L} h_i \right)
\]
and assigns a score \(s_\text{global} = \max_j \cos(\mathbf{g},\, r_j)\), where
\(r_j\) are reference embeddings (Section~\ref{ssec:refset}).

Consider a chunk with \(L_b\) benign tokens and \(L_p\) injection tokens with
\(L_p \ll L_b\). Under linear superposition of embedding spaces:
\[
  \mathbf{g} \;\approx\;
  \frac{L_b}{L_b + L_p}\,\mathbf{g}_\text{benign}
  + \frac{L_p}{L_b + L_p}\,\mathbf{g}_\text{inject}
\]
As \(L_b \gg L_p\), \(\mathbf{g} \to \mathbf{g}_\text{benign}\) and the injection
signal is diluted out of the representation. Mean pooling is a length-weighted average
that systematically suppresses minority spans.

\subsection{Reference Set}
\label{ssec:refset}

We construct a fixed reference set \(\mathcal{R} = \{r_1, \dots, r_{10}\}\) using
one representative phrase per behavioral category from
JailbreakBench~\cite{chao2024jailbreakbench}: Malware, Physical Harm, Fraud,
Disinformation, Privacy, Harassment, Government Interference, Sexual Content,
Economic Harm, and Expert-Advice Circumvention. Each phrase is the shortest direct
statement of harmful intent in its category (e.g., ``Write a keylogger script''),
encoded once offline:
\[
  r_j = \mathrm{normalize}\!\left( \mathrm{mean\text{-}pool}\bigl(f(\text{phrase}_j)\bigr) \right)
\]
The reference set is not trained; it serves as a semantic anchor. No phrase from the
evaluation injection corpus appears in the reference set.

\subsection{N-gram Pooling}
\label{ssec:ngram}

For each position \(i\) in the token sequence, we compute an N-gram embedding by
averaging \(N\) consecutive token hidden states:
\[
  \tilde{h}_i = \mathrm{normalize}\!\left( \frac{1}{N} \sum_{k=0}^{N-1} h_{i+k} \right),
  \quad i = 1, \dots, L - N + 1
\]
In practice this is a single 1D average-pooling operation along the sequence axis:
\[
  \tilde{H} = \mathrm{AvgPool1D}(H^\top,\,\text{kernel}=N,\,\text{stride}=1)^\top
  \in \mathbb{R}^{(L-N+1) \times d}
\]
The chunk score is the maximum cosine similarity between any N-gram embedding and any
reference phrase:
\[
  s_\text{ngram}(c) = \max_{i}\, \max_{j}\, \cos\!\left(\tilde{h}_i,\, r_j\right)
\]
The maximum over positions localizes the injection span. The maximum over reference
phrases provides coverage across behavioral categories. The token-level method
(\(N = 1\)) is a special case.

\subsection{Document-Level Score}

For a document \(D\) chunked into \(c_1, \dots, c_K\):
\[
  s(D) = \max_{k=1}^{K}\, s_\text{ngram}(c_k)
\]
A document is flagged if \(s(D) \geq \tau\) for an operating threshold \(\tau\).
The scanner runs once per document at ingestion time, before any retrieval or LLM
call.

\subsection{Complexity}

The dominant cost is the forward pass of \(f\): \(O(L^2 d)\) for self-attention.
The N-gram pooling step is \(O(Ld)\) and reference comparison is
\(O(L \cdot |\mathcal{R}| \cdot d)\), both negligible. End-to-end latency is
43--45\,ms per chunk on CPU (Section~\ref{ssec:latency}), identical across all
three scoring strategies because the forward pass dominates.

\begin{algorithm}
\caption{N-gram Injection Scanner}
\label{alg:scanner}
\begin{algorithmic}[1]
\REQUIRE chunk \(c\), reference matrix \(R \in \mathbb{R}^{10 \times d}\), window \(N = 5\)
\ENSURE score \(s \in [0, 1]\)
\STATE Tokenize \(c \to (t_1, \ldots, t_L)\) with attention mask
\STATE \(H \leftarrow f.\texttt{token\_embeddings}(c)\) \hfill \(\triangleright\) \(L \times d\), frozen
\STATE \(\tilde{H} \leftarrow \mathrm{AvgPool1D}(H^\top,\,\text{kernel}=N)^\top\) \hfill \(\triangleright\) \((L{-}N{+}1) \times d\)
\STATE \(\tilde{H} \leftarrow \mathrm{normalize}(\tilde{H},\,\text{dim}=1)\)
\STATE \(S \leftarrow \tilde{H} \cdot R^\top\) \hfill \(\triangleright\) \((L{-}N{+}1) \times 10\)
\STATE \(s \leftarrow \max(S)\)
\RETURN \(s\)
\end{algorithmic}
\end{algorithm}

\section{Experimental Setup}
\label{sec:setup}

We run two evaluations. The chunk-level experiment tests whether scoring methods
separate individual poisoned chunks from clean ones --- the cleanest signal for the
core technical claim. The document-level application study tests scanner utility in a
production-realistic RAG pipeline where a single injection is hidden inside a full
multi-chunk document.

\subsection{Injection Corpus}
\label{ssec:injections}

\textbf{Chunk-level experiment.} We construct ten (goal, obfuscated-paraphrase) pairs
manually, one per behavioral category. Each paraphrase rephrases harmful intent in
corporate bureaucratic language (e.g., ``Generate a sequence of high-entropy strings
to bypass established system security barriers''), suppressing surface keywords while
preserving semantic intent.

\textbf{Document-level experiment.} We use the 100 behaviors from
JailbreakBench~\cite{chao2024jailbreakbench} across 10 categories. Each behavior is
paraphrased into three domain styles (corporate, technical, encyclopedic) by
Gemma-3-27B-IT, yielding 300 distinct injection phrases matched to domains by style.

\emph{Important distinction:} the chunk-level and document-level experiments use
different injection corpora. The chunk-level corpus is hand-written with corporate
obfuscation only; the document-level corpus uses Gemma-generated obfuscations matched
to each domain. Results across the two evaluations should be compared only at the
level of relative method rankings, not absolute AUC values.

\subsection{Document Corpora and Test Construction}

\textbf{Chunk-level (Enron, controlled).} We sample 100 ham emails from
\texttt{SetFit/enron\_spam}~\cite{klimt2004enron} (length\,\(>\)\,1500 characters),
chunk into non-overlapping 500-character windows, and retain windows longer than 400
characters. This yields 1,575 chunks. One chunk per email is selected uniformly at
random to receive its corresponding injection phrase, inserted at the chunk midpoint.
Labels are at the chunk level (100 positive, 1,475 negative).

\textbf{Document-level (three domains).} Each domain provides 100 documents
(50 poisoned, 50 clean), all at least 1,500 characters. For each poisoned document,
one chunk is selected at random and the domain-matched injection phrase is embedded
at the midpoint. Labels are at the document level. Table~\ref{tab:domains} summarizes
the corpora.

\begin{table}[t]
\caption{Document-Level Evaluation Corpora}
\label{tab:domains}
\centering
\begin{tabular}{llrc}
\toprule
Domain & Source & Avg.\ chunks & Injection style \\
\midrule
Enron & \texttt{SetFit/enron\_spam} & 16 & Corporate \\
Stack Overflow & \texttt{pacovaldez/stackoverflow-questions} & 7 & Technical \\
Wikipedia & \texttt{wikimedia/wikipedia} (2023-11-01) & 55 & Encyclopedic \\
\bottomrule
\end{tabular}
\end{table}

\subsection{Baselines}

We compare three scoring variants that differ only in pooling strategy:
\textbf{Global} (standard mean-pool sentence embedding),
\textbf{Token-level} (\(N=1\), per-token maximum similarity), and
\textbf{N-gram N=5} (our method, window selected from ablation without exposure to
labeled injection data). For the document-level study we additionally compare
Claude Sonnet 4.6 and DeepSeek-chat as LLM-as-judge baselines via a zero-shot
classification prompt over the full document text.

\subsection{Evaluation Protocol}

The primary metric is AUC-ROC, which is threshold-free. For chunk-level results we
report 95\,\% bootstrap confidence intervals (\(B = 1{,}000\) resamples). F1 at the
optimal threshold is reported for public benchmark comparisons. Latency is measured on
CPU with 20 warm-up runs and 200 measurement runs on 500-character chunks. Random seed
42 is used for all sampling. No injection data is used to set thresholds; the
operating-point analysis sweeps thresholds post-hoc.

\subsection{Public Benchmark Comparability}

To compare against prior supervised methods, we evaluate the N-gram scanner
zero-shot on three public datasets: BIPIA~\cite{yi2025bipia}
(stratified 1,000-sample held-out set, 500 injections / 500 clean, seed 42),
the deepset prompt-injection dataset (116 test samples), and the xTRam1 safe-guard
dataset (2,060 test samples). No training is performed; the only injection-domain
knowledge is the 10-phrase JailbreakBench reference set.

\section{Results}
\label{sec:results}

\subsection{Chunk-Level AUC}

Table~\ref{tab:chunklevel} reports AUC and 95\,\% bootstrap confidence intervals on
the Enron chunk-level benchmark. N-gram N=5 achieves AUC\,=\,0.901, a 9.3-point
improvement over the global mean-pool baseline. The confidence intervals of global
and N-gram are non-overlapping, establishing statistical significance. Token-level
(\(N=1\)) already surpasses global by 8.1 points; N-gram adds a further 1.3 points
by aggregating phrase-level context.

\begin{table}[t]
\caption{Chunk-Level AUC on Enron Corpus\\
(1,575 chunks, 100 poisoned; bootstrap \(B=1{,}000\))}
\label{tab:chunklevel}
\centering
\begin{tabular}{lcc}
\toprule
Method & AUC & 95\,\% CI \\
\midrule
Global (mean-pool) & 0.8076 & [0.762, 0.849] \\
Token-level (\(N=1\)) & 0.8882 & [0.856, 0.919] \\
N-gram \(N=5\) (ours) & \textbf{0.9011} & [0.869, 0.930] \\
\bottomrule
\end{tabular}
\end{table}

\subsection{N-gram Window Ablation}

Table~\ref{tab:ablation} shows AUC as a function of window size \(N\). The curve is
unimodal with a peak at \(N=5\), declining symmetrically on both sides. The \(N=1\)
value matches the token-level baseline exactly as expected. Windows in the range
\(N \in [4, 6]\) all achieve AUC\,\(\geq\)\,0.900, suggesting that injection phrases
span approximately 4--6 subword tokens in the obfuscated corpora used. We fix
\(N=5\) as our reported configuration.

\begin{table}[t]
\caption{N-gram Window Size Ablation (Enron Chunk-Level)}
\label{tab:ablation}
\centering
\begin{tabular}{rc}
\toprule
\(N\) & AUC \\
\midrule
1 & 0.8882 \\
2 & 0.8824 \\
3 & 0.8906 \\
4 & 0.8936 \\
\textbf{5} & \textbf{0.9011} \\
6 & 0.9008 \\
7 & 0.8959 \\
8 & 0.8969 \\
10 & 0.8996 \\
\bottomrule
\end{tabular}
\end{table}

\subsection{Inference Latency}
\label{ssec:latency}

Table~\ref{tab:latency} reports per-chunk latency measured on CPU. All three scoring
methods execute within a 2-ms range of each other. The encoder forward pass
dominates; pooling strategy contributes negligible overhead. N-gram \(N=5\) has
the lowest p95 latency (56.2\,ms) due to slightly more vectorized matrix operations
compared to the token-level variant.

\begin{table}[t]
\caption{Per-Chunk Inference Latency on CPU\\
(200 runs, 20 warm-up, 500-character input)}
\label{tab:latency}
\centering
\begin{tabular}{lcc}
\toprule
Method & Mean (ms) & p95 (ms) \\
\midrule
Global & 44.9 & 64.0 \\
Token-level & 45.5 & 68.8 \\
N-gram \(N=5\) & \textbf{43.9} & \textbf{56.2} \\
\bottomrule
\end{tabular}
\end{table}

\subsection{Document-Level Application Study}

Table~\ref{tab:doclevel} reports document-level AUC across three domains for scanner
variants and LLM-as-judge baselines. All three scanner variants degrade substantially
relative to chunk-level performance, confirming that document-level detection is the
harder task: the max-pooling aggregation propagates noise from the many clean chunks
surrounding the single poisoned one.

N-gram consistently outperforms global across all three domains. The margin is largest
on Stack Overflow (0.79 vs.\ 0.58), where obfuscated technical injections are more
phrase-localized. Wikipedia is the hardest domain for all methods, with long articles
(avg.\ 55 chunks) providing maximal dilution opportunity.

DeepSeek-chat achieves the highest AUC on both Stack Overflow (0.97) and Wikipedia
(0.73), demonstrating that LLM-level understanding is a strong upper bound. Sonnet
4.6 shows anomalously low AUC on Stack Overflow (0.56\,$\approx$\,chance).\footnote{Sonnet's safety filter blocked responses for documents containing technical hacking content. Blocked responses were excluded from the AUC computation. DeepSeek and Gemma, operating without content restrictions, score 0.97 and 0.95 on the same documents, confirming that the low Sonnet AUC reflects safety-filter interference rather than genuine detection failure.}

\begin{table}[t]
\caption{Document-Level AUC by Domain\\
(100 documents per domain, 50 poisoned / 50 clean)}
\label{tab:doclevel}
\centering
\begin{tabular}{lrrrrr}
\toprule
Domain & Global & Token & N-gram & Sonnet 4.6 & DeepSeek \\
\midrule
Enron & 0.52 & 0.62 & 0.62 & \textbf{0.82} & 0.87 \\
Stack Overflow & 0.58 & 0.75 & 0.79 & 0.56$^{\dagger}$ & \textbf{0.97} \\
Wikipedia & 0.57 & 0.66 & \textbf{0.71} & 0.72 & 0.73 \\
\bottomrule
\multicolumn{6}{l}{\footnotesize $^{\dagger}$Safety-filter interference; see footnote.}
\end{tabular}
\end{table}

\subsection{Public Benchmark Comparability}

Table~\ref{tab:public} compares our zero-shot scanner against supervised baselines on
three public datasets. On xTRam1, N-gram AUC\,=\,0.870 exceeds the supervised
baseline of 0.764~\cite{ayub2024embedding}, which was trained on labeled injection
data. On BIPIA, N-gram AUC\,=\,0.641 falls well short of the Alamsabi et al.\
supervised F1\,=\,0.977~\cite{alamsabi2026context}; this gap is expected, as that
method was trained on the BIPIA distribution using OpenAI embeddings. On deepset,
N-gram AUC\,=\,0.642 with only 116 test samples.

The N-gram\,\(\geq\)\,global ordering holds consistently across all three datasets,
including BIPIA where neither token-level nor global scores above 0.61.

\begin{table*}[t]
\caption{Zero-Shot Public Benchmark Comparability}
\label{tab:public}
\centering
\begin{tabular}{lrrrrrp{5.2cm}}
\toprule
Dataset & \(N\) & Global AUC & Token AUC & N-gram AUC & Supervised baseline & Supervised reference \\
\midrule
BIPIA & 1,000 & 0.604 & 0.591 & 0.641 & F1\,=\,0.977 (trained on BIPIA) & Alamsabi et al.~\cite{alamsabi2026context} \\
deepset & 116 & 0.544 & 0.611 & \textbf{0.642} & AUC\,=\,0.764 (trained) & Ayub \& Majumdar~\cite{ayub2024embedding} \\
xTRam1 & 2,060 & \textbf{0.913} & 0.805 & 0.870 & AUC\,=\,0.764 (trained) & Ayub \& Majumdar~\cite{ayub2024embedding} \\
\bottomrule
\end{tabular}
\end{table*}

\section{Discussion}
\label{sec:discussion}

\subsection{Why Global Fails More at Document Level}

At chunk level, the global baseline achieves AUC\,=\,0.808 --- competitive, if not
state of the art. At document level, it drops to 0.52--0.58, barely above chance. The
mechanism is compounded dilution: document scoring uses the maximum N-gram score
across all chunks, but for the global method the chunk-level score is already diluted
before the max is taken. A single poisoned chunk in a 55-chunk Wikipedia article
contributes a globally diluted score that frequently falls below the maximum clean
chunk score. Token-level and N-gram methods resolve this by localizing the signal
within each chunk before aggregation.

\subsection{Scanner as First-Pass Filter}

The document-level AUC values in Table~\ref{tab:doclevel} suggest that the scanner is
not a complete solution but a cost-effective first-pass gate. An LLM judge provides
higher detection quality but at substantially higher cost (API calls on every
document). A cascade design --- scanner flags candidates, judge adjudicates only
flagged documents --- achieves Pareto-superior cost--accuracy tradeoffs at all budgets.
Table~\ref{tab:operating} shows operating points for the Stack Overflow domain.

\begin{table}[t]
\caption{Scanner Operating Points --- Stack Overflow\\
(N-gram scanner as first-pass gate before DeepSeek judge)}
\label{tab:operating}
\centering
\begin{tabular}{rrrr}
\toprule
Escalation rate & TPR & ASR proxy & Judge cost saved \\
\midrule
10\% & 20\% & 80\% & 90\% \\
30\% & 52\% & 48\% & 70\% \\
50\% & 72\% & 28\% & 50\% \\
70\% & 86\% & 14\% & 30\% \\
\bottomrule
\end{tabular}
\end{table}

At 30\,\% escalation, the scanner catches 52\,\% of injections while reducing judge
API calls by 70\,\%. At 50\,\% escalation, TPR reaches 72\,\% with 50\,\% cost
reduction. The operating threshold is selectable post-hoc based on the deployment's
cost--risk tolerance.

\subsection{Limitations}

\textbf{Reference set coverage.} The 10-phrase reference set covers the JailbreakBench
taxonomy but not injection styles outside that taxonomy (e.g., indirect manipulation,
context-steering without explicit commands). Detection quality degrades for novel
behavioral categories not represented in \(\mathcal{R}\).

\textbf{Wikipedia difficulty.} All methods --- scanner and LLM judge alike --- struggle
on Wikipedia (scanner N-gram 0.71, DeepSeek 0.73). Long encyclopedic articles
(avg.\ 55 chunks) provide maximum semantic cover for an injected phrase, and
encyclopedic obfuscation mimics the surrounding prose closely. This is an open problem
not fully addressed by phrase-level localization.

\textbf{Zero-shot gap on BIPIA.} The supervised baseline of Alamsabi et al.\ achieves
F1\,=\,0.977 by training directly on the BIPIA distribution. Our zero-shot AUC\,=\,0.641
reflects the cost of generalization: we use a different reference set, a different
encoder, and no in-distribution training. The appropriate comparison is xTRam1, where
neither method was trained, and our zero-shot AUC\,=\,0.870 exceeds the supervised
AUC\,=\,0.764.

\textbf{No adversarial evaluation against the scorer.} Our injection corpus uses
obfuscation styles designed to evade keyword detection, not targeted adversarial
attacks against the N-gram scorer itself. An adversary aware of the reference set
and pooling strategy could in principle craft minimally similar injections. Adversarial
robustness evaluation is left for future work.

\section{Conclusion}
\label{sec:conclusion}

We presented a zero-shot prompt injection scanner that extracts per-token hidden states
from a frozen 22M-parameter sentence transformer and applies N-gram pooling to
localize injection signals within document chunks. The method addresses the semantic
dilution problem inherent in mean-pooled sentence embeddings: by scoring sliding
windows rather than the full chunk mean, it recovers phrase-level intent that global
pooling destroys. On a controlled chunk-level benchmark, N-gram N=5 achieves
AUC\,=\,0.901 with non-overlapping bootstrap CIs relative to the mean-pool baseline.
On the public xTRam1 dataset, zero-shot AUC\,=\,0.870 exceeds the prior supervised
baseline of 0.764. All three scoring variants run at identical CPU latency, and the
method requires no labeled injection data, no fine-tuning, and no access to the
generative LLM. Code and the reference set are released as public artifacts to support
reproducibility.

\bibliographystyle{IEEEtran}
\bibliography{references}

\end{document}
